\chapter{Modeling}

\section{Red Flag Detection}
(Sezione in fase di stesura...)

\section{Depression Prediction}

\subsection{User-Based Splitting e Baseline Regressiva}

La costruzione dei modelli predittivi per il modulo \textit{Depression Prediction} ha richiesto un approccio metodologico rigoroso per gestire la natura intrinsecamente longitudinale dei dati fisiologici ed emotivi. Poiché i check-in degli utenti costituiscono serie temporali, un partizionamento ingenuo riga per riga porterebbe a un massiccio \textit{data leakage}, in cui il modello apprenderebbe l'identità del paziente anziché i pattern patologici della depressione.

Per neutralizzare questo rischio alla radice, il workflow di SINTON-IA ha previsto l'aggregazione di tutte le misurazioni temporali di ciascun individuo in un unico vettore di feature riassuntive (eseguito in fase di \textit{Feature Engineering}). Di conseguenza, il dataset utilizzato per il modeling presenta una struttura in cui \textbf{ogni record corrisponde a un singolo paziente unico}. In questa configurazione, lo split tra le partizioni di training e test diventa \textbf{User-Based per costruzione}: garantendo l'atomicità del dato a livello utente, si assicura che nessuna informazione riguardante un singolo individuo possa essere contemporaneamente presente in entrambi i set, preservando l'integrità della valutazione e la capacità di generalizzazione del sistema.

\subsubsection*{Dimensioni delle Partizioni}
I pazienti sono stati estratti dal dataset preparato e suddivisi in due set distinti. Seguendo le best practices di letteratura per dataset di queste dimensioni, abbiamo optato per una ripartizione stratificata:
\begin{itemize}
    \item \textbf{Training Set (80\%)}: $80.000$ record (utilizzato per lo sviluppo e l'ottimizzazione).
    \item \textbf{Test Set (20\%)}: $20.000$ record (congelato e impiegato esclusivamente per la valutazione finale).
\end{itemize}

\subsubsection*{Distribuzione del Target}
Ultimato lo split, è stato verificato che la distribuzione della variabile target (il punteggio \textbf{PHQ-9 Total}) rimanesse inalterata in entrambe le partizioni. La Figura~\ref{fig:dp_target_dist} riporta le funzioni di densità stimate (KDE) del target, confermando l'assenza di bias indotti dal processo di separazione dei pazienti.

\begin{figure}[H]
    \centering
    \includegraphics[width=0.8\textwidth]{dp_target_distribution.png}
    \caption{Distribuzione della variabile target PHQ-9 nelle partizioni di Training e Test.}
    \label{fig:dp_target_dist}
\end{figure}

\subsubsection*{Baseline Regressiva Interpretabile}
Prima di addestrare algoritmi ad elevata complessità, è stata definita una \textbf{Baseline} di riferimento basata sulla \textbf{Ridge Regression}. La Ridge Regression è un modello lineare "regolarizzato": a differenza della regressione standard, essa introduce una piccola penalità numerica sui coefficienti del modello. Questo accorgimento serve a mantenere le predizioni stabili e a evitare l'overfitting (il fenomeno per cui il modello memorizza i dati di training invece di imparare a generalizzare), rendendolo un benchmark solido e affidabile.

La valutazione della baseline è stata condotta tramite \textbf{KFold Cross-Validation} a 5 split sul Training Set. Questa tecnica consiste nel dividere i dati in 5 gruppi ("fold") e addestrare il modello 5 volte: ogni volta, 4 gruppi vengono usati per l'apprendimento e il rimanente per la verifica, a rotazione. Questo processo "round-robin" garantisce che ogni record venga testato almeno una volta, fornendo una stima della performance molto più robusta rispetto a un singolo test casuale.

I risultati medi della cross-validation per il modello baseline sono stati i seguenti:
\begin{itemize}
    \item \textbf{Mean Absolute Error (MAE)}: $2.6793 \pm 0.0096$. Il MAE rappresenta l'errore medio "puro" (la distanza media tra valore predetto e reale), espresso nella stessa unità di misura del PHQ-9. Un valore di circa 2.6 significa che, mediamente, il modello sbaglia di meno di 3 punti sulla scala diagnostica.
    \item \textbf{Root Mean Squared Error (RMSE)}: $3.3861 \pm 0.0079$. L'RMSE è simile al MAE ma penalizza più duramente gli errori grandi. Essendo leggermente superiore al MAE, ci indica che il modello commette saltuariamente errori più marcati, sebbene rimanga mediamente preciso.
    \item \textbf{Coefficiente di Determinazione ($R^2$)}: È un indice che ci dice quanto il modello sia più "intelligente" rispetto a una semplice media statistica. Se per ogni nuovo paziente ci limitassimo a "tirare a indovinare" usando sempre il punteggio medio di tutti gli altri pazienti, avremmo un $R^2$ pari a 0. Il valore di $0.6171$ indica che la nostra baseline lineare aggiunge un valore significativo, riuscendo a spiegare circa il 62\% della variabilità dei punteggi reali.
\end{itemize}

Queste misurazioni costituiscono il punto di partenza per quantificare i miglioramenti apportati da architetture più avanzate.

\subsection{Selezione del Modello e Hyperparameter Tuning}

Successivamente all'istituzione della baseline lineare, la sperimentazione ha previsto l'addestramento e il confronto di modelli \textit{ensemble} basati su alberi decisionali, tipicamente più adatti a catturare relazioni non lineari:
\begin{itemize}
    \item \textbf{Random Forest}: Opera come un "comitato" di centinaia di alberi decisionali indipendenti. Ogni albero formula una predizione e il risultato finale è la media di tutte; questo riduce drasticamente l'errore casuale.
    \item \textbf{LightGBM}: A differenza del Random Forest, gli alberi non sono indipendenti ma vengono creati in sequenza. Ogni nuovo albero cerca di correggere gli errori commessi dai precedenti. È un approccio molto più mirato e computazionalmente efficiente.
\end{itemize}

\subsubsection*{Confronto Iniziale sul Training Set}
I modelli candidati sono stati valutati con la stessa configurazione di KFold CV. La Tabella~\ref{tab:model_comparison} riassume le performance medie. Come atteso, i metodi ensemble hanno superato nettamente la regressione lineare.

\begin{table}[H]
    \centering
    \begin{tabular}{|l|c|c|}
        \hline
        \textbf{Modello Candidato} & \textbf{MAE (M $\pm$ SD)} & \textbf{RMSE (M $\pm$ SD)} \\
        \hline
        Ridge Regression (Baseline) & $2.6793 \pm 0.0096$ & $3.3861 \pm 0.0079$ \\
        Random Forest               & $2.2509 \pm 0.0158$ & $2.8906 \pm 0.0181$ \\
        LightGBM                    & $2.3120 \pm 0.0180$ & $2.9560 \pm 0.0200$ \\
        \hline
    \end{tabular}
    \caption{Confronto delle prestazioni medie in CV (KFold) sul Training Set.}
    \label{tab:model_comparison}
\end{table}

\subsubsection*{Ottimizzazione del LightGBM e Sample Weights}
Sebbene il Random Forest abbia mostrato metriche iniziali leggermente superiori, il \textbf{LightGBM Regressor} è stato selezionato come modello definitivo per SINTON-IA per diverse ragioni strategiche. Innanzitutto, il LightGBM gestisce nativamente i valori mancanti, una caratteristica frequente nei dati reali di telemetria clinica. In secondo luogo, il suo algoritmo di crescita "leaf-wise" (che espande i nodi più "promettenti" invece di quelli più vicini alla radice) è statisticamente più incline a catturare pattern complessi e sottili tipici del comportamento umano, non appena viene rifinito tramite ottimizzazione.

La sintonia fine degli iperparametri è stata delegata ad \textbf{Optuna}, un framework di ottimizzazione bayesiana. Diversamente dai classici tentativi casuali, Optuna agisce in modo "intelligente": analizza i risultati dei test precedenti per restringere lo spazio di ricerca verso le configurazioni più promettenti. Questo approccio permette di trovare il miglior assetto del modello senza sprecare risorse computazionali, imparando letteralmente dai propri "errori" di configurazione. Sono stati eseguiti \textbf{50 cicli (trials)} di ottimizzazione, valutati tramite una validazione K-Fold a 3 split.

Infine, per allineare il modello alle stringenti necessità cliniche di SINTON-IA, abbiamo introdotto i \textbf{Sample Weights}. Ai pazienti con depressione grave (score $\ge 15$) è stato assegnato un peso raddoppiato durante l'addestramento. La decisione ha una forte motivazione clinica: l'obiettivo non è solo essere precisi "in generale", ma soprattutto evitare i falsi negativi nei soggetti critici. Un errore di sottovalutazione per un paziente grave è considerato molto più grave di un errore per un paziente sano; raddoppiando il peso di questi ultimi, costringiamo matematicamente l'algoritmo a dare priorità assoluta alla correttezza diagnostica nelle fasce d'emergenza.

I parametri finali estratti dal tuning bayesiano definiscono un modello profondamente ottimizzato. L'assetto definitivo è stato addestrato integralmente sull'intero Training Set, pesato per gravità clinica.

Un vantaggio significativo dei moderni modelli Tree-Based è la possibilità di estrarre e quantificare rigorosamente l'importanza di ciascuna variabile predittiva. Al posto delle metriche standard (che spesso sovrastimano variabili con codici univoci), è stato integrato il framework \textbf{SHAP} (\textit{SHapley Additive exPlanations}), basato sulla Teoria dei Giochi (Figura~\ref{fig:dp_feature_importance}).

\begin{figure}[H]
    \centering
    \includegraphics[width=0.8\textwidth]{dp_feature_importance.png}
    \caption{SHAP Summary Plot per le dinamiche comportamentali nel LightGBM. L'asse X indica l'impatto direzionale sullo score regressivo diagnostico.}
    \label{fig:dp_feature_importance}
\end{figure}

Il concetto alla base di SHAP è quello di considerare ogni variabile come un "giocatore" di una squadra (il modello) che collabora per ottenere un risultato (la predizione). SHAP calcola matematicamente quanto "merito" spetta a ciascuna variabile per aver spostato il punteggio finale rispetto alla media. Questo garantisce una trasparenza senza precedenti: non sappiamo solo quali feature sono importanti, ma anche se aumentano o diminuiscono l'allarme depressivo. 

Il Summary Plot di SHAP ci rivela \textit{come} varia l'allarme depressione rispetto ai valori assunti dalle feature. Il tono medio (\texttt{valence\_mean}) si conferma il fattore statisticamente più influente: valori bassi di valenza emozionale (toni scuri nel grafico) spingono la predizione verso punteggi PHQ-9 più alti. Di pari rilevanza clinica è il peso attribuito all'assenza di dati (\texttt{missing\_ratio}) e all'ampiezza degli sbalzi d'umore (\texttt{intensity\_mean}). Tali risultati dimostrano esplicitamente come SHAP conferisca interpretabilità direzionale: un innalzamento patologico di queste misurazioni agisce da precursore critico del deterioramento clinico, venendo riconosciuto dal sistema come primario fattore prognostico di scompenso emotivo.

\subsection{Valutazione Finale e Analisi dei Residui}
Il modello LightGBM, pesato e ottimizzato bayesianamente, è stato infine validato sul Test Set indipendente per valutarne le reali capacità di generalizzazione e la prontezza diagnostica per la messa in produzione.

\subsubsection*{Performance Generali}
I risultati sul Test Set (Tabella~\ref{tab:final_evaluation_metrics}) confermano l'affidabilità del framework. Con un $R^2$ incrementato e ottimi controlli su MAE e RMSE, il sistema si rivela in grado di spiegare larga parte della varianza associata allo score depressivo avvalendosi esclusivamente di metadati comportamentali (il diario emozionale) prima ancora che il paziente compili i test formali.

\begin{table}[H]
    \centering
    \begin{tabular}{|l|c|c|c|}
        \hline
        \textbf{Configurazione} & \textbf{MAE} & \textbf{RMSE} & \textbf{$R^2$} \\
        \hline
        Baseline (Ridge CV) & 2.6793 & 3.3861 & 0.6171 \\
        \textbf{LightGBM Ottimizzato (Test Set)} & \textbf{2.2098} & \textbf{2.8273} & \textbf{0.7321} \\
        \hline
        \textit{Miglioramento Assoluto} & \textbf{-0.4695} & \textbf{-0.5588} & \textbf{+0.1150} \\
        \hline
    \end{tabular}
    \caption{Prestazioni in inferenza del LightGBM definitivo sul Test Set.}
    \label{tab:final_evaluation_metrics}
\end{table}

\subsubsection*{Confronto con le Soglie di Accettabilità}
I risultati ottenuti testimoniano che il modello LightGBM soddisfa e supera ampiamente le soglie di accettabilità per la validazione diagnostica postulate nel contesto analitico (Capitolo 2). A fronte dei requisiti di tolleranza di $\text{MAE} \le 3.0$ e $\text{RMSE} \le 4.5$, il sistema definitivo in inferenza registra performance nettamente superiori. Tali margini prestazionali non solo confermano la fattibilità ingegneristica dell'approccio, ma ne legittimano concretamente l'impiego operativo entro i flussi di SINTONIA. Il modello è adeguatamente maturo per erogare previsioni stabili, supportando in maniera non invasiva e precauzionale le decisioni cliniche senza inondare i professionisti sanitari di falsi allarmi lievi derivanti da un fitting grossolano.

\subsubsection*{Distribuzione dell'Errore e Assenza di Bias Sistematica}
L'analisi dei residui (la differenza tra ciò che il modello predice e il valore clinico reale) è fondamentale per capire se l'algoritmo ha delle "preferenze" o se è equilibrato. L'istogramma degli errori sul Test Set (Figura~\ref{fig:dp_residual_dist}) mostra una distribuzione a campana centrata sullo zero. Questo significa che il modello è sprovvisto di quello che in gergo tecnico si chiama \textit{bias}: non tende né a sovrastimare né a sottostimare sistematicamente, ma distribuisce i suoi piccoli errori in modo equo. Immaginando una bilancia, il modello è perfettamente in equilibrio; quando sbaglia, lo fa in modo casuale e contenuto, senza mai "pendere" troppo da una parte o dall'altra.

\begin{figure}[H]
    \centering
    \begin{minipage}{0.48\textwidth}
        \centering
        \includegraphics[width=\linewidth]{dp_residual_distribution.png}
        \caption{Istogramma dei residui ($y_{true} - y_{pred}$) valutato sul Test Set.}
        \label{fig:dp_residual_dist}
    \end{minipage}\hfill
    \begin{minipage}{0.48\textwidth}
        \centering
        \includegraphics[width=\linewidth]{dp_residual_plot.png}
        \caption{Residual Plot e dispersione locale in base al target predetto.}
        \label{fig:dp_residual_plot}
    \end{minipage}
\end{figure}

Il Residual Plot in Figura~\ref{fig:dp_residual_plot} testimonia una \textit{omochedasticità} ragionevole: la varianza dell'errore resta sufficientemente costante lungo i vari indici di severità, rassicurando l'affidabilità delle soglie generate dai \textit{nudges} futuri.

\subsubsection*{Analisi Stratificata per Fascia PHQ-9 e Implicazioni Cliniche}
Per capire l'utilità reale del sistema in un ospedale o in una clinica, non basta guardare la precisione media: dobbiamo sapere come si comporta il modello per ogni tipologia di paziente. Analizzare i risultati per ogni fascia di gravità (lieve, moderata, grave) è fondamentale perché un errore di 2 punti ha un peso clinico molto diverso se il paziente è sano (punteggio basso) o se è già in una zona di criticità (punteggio alto). L'obiettivo di SINTON-IA non è solo misurare l'umore, ma fungere da "salvagente" per i casi più complessi.

\begin{table}[H]
    \centering
    \begin{tabular}{|l|c|c|}
        \hline
        \textbf{Fascia di Severità PHQ-9 (Reale)} & \textbf{Numero Pazienti} & \textbf{Mean Absolute Error} \\
        \hline
        Minima (0-4) & 11.038 & $1.7661$ \\
        Lieve (5-9) & 4.987 & $2.4967$ \\
        Moderata (10-14) & 2.372 & $3.4666$ \\
        Moderatamente Grave (15-19) & 1.010 & $2.3585$ \\
        Grave (20-27) & 593 & $2.7739$ \\
        \hline
    \end{tabular}
    \caption{Disaggregazione dell'Errore Medio Assoluto nelle bande nosologiche PHQ-9.}
    \label{tab:mae_by_band}
\end{table}

Come si evince dalla Tabella~\ref{tab:mae_by_band}, l'inserimento dei class weights ha prodotto una riduzione significativa dei falsi sani: per i pazienti nella fascia \textbf{Moderatamente Grave} e \textbf{Grave}, il modello mantiene una reattività diagnostica proattiva (errori compresi tra 2.5 e 3 punti di scarto della scala), proteggendo questi casi dall'essere banalizzati o etichettati erroneamente in fasce non d'intervento.

\begin{figure}[H]
    \centering
    \includegraphics[width=0.85\textwidth]{dp_scatter_pred_vs_true.png}
    \caption{Scatterplot di predizioni finali sovrapposte alle fasce cliniche PHQ-9.}
    \label{fig:dp_scatter_bands}
\end{figure}

\subsection{Limiti dell'Approccio e Sviluppi Futuri}
L'architettura sviluppata per la \textit{Depression Prediction} dimostra la fattibilità tecnica di inferire punteggi clinici complessi analizzando esclusivamente tracce comportamentali asincrone. Ciononostante, si evidenziano specifici limiti ingegneristici che dovranno essere affrontati in vista di futuri aggiornamenti della piattaforma: 

\begin{itemize}
    \item \textbf{Generazione Sintetica dei Dati}: A causa dello stadio iniziale del progetto SINTONIA (start-up), il modello è stato validato mediante simulazioni digitali (\textit{Digital Twins}), ancorché accuratissime dal punto di vista probabilistico. Ne consegue un limite epistemologico intrinseco: l'Intelligenza Artificiale non ha ancora valutato la stocasticità caotica o le incongruenze umane (imprecisioni logiche, typo, o variazioni socio-anagrafiche complesse) che insorgeranno nel lancio in via sperimentale nell'ASL. Al varo clinico ufficiale, il LightGBM dovrà obbligatoriamente subire una fase di massiccio \textit{re-training} sui primi flussi di telemetria reali per compensare il possibile data shift.
    \item \textbf{Gestione del Cold Start}: Il modello odierno vanta elevata predittività pre-allertando i medici, ma vincola la propria analisi ad almeno due settimane ($\sim$14 giorni) di tracciatura dati per estrarre features temporali stabili come deviazione standard o tendenze in regressione polinomiale. Un paziente vulnerabile neoiscritto (\textit{Cold Start}) rimarrà scoperto per diversi giorni dal monitoraggio algoritmico puro, affidando giocoforza le sue segnalazioni d'acuzie al comparto parallelo di \textit{Red Flag Detection} sul testo lineare. 
\end{itemize}
L'integrazione parallela delle features linguistiche grezze con le derivate temporali nel medesimo regressore resta la sfida scientifica cardine prospettata per i prossimi avanzamenti di prodotto.
